\documentclass[10pt,a4paper]{amsart}
\usepackage{amsmath,amssymb,amsthm,graphicx}
% For commutative diagrams we use Xy-pic
\usepackage[all]{xy}

% We use 2cell for 2-commutative diagrams.
\xyoption{2cell}
% Theorem environments.
%
\theoremstyle{plain}
\newtheorem{theorem}[subsection]{Theorem}
\newtheorem{proposition}[subsection]{Proposition}
\newtheorem{lemma}[subsection]{Lemma}

\theoremstyle{definition}
\newtheorem{definition}[subsection]{Definition}
\newtheorem{example}[subsection]{Example}
\newtheorem{exercise}[subsection]{Exercise}
\newtheorem{situation}[subsection]{Situation}

\theoremstyle{remark}
\newtheorem{remark}[subsection]{Remark}
\newtheorem{remarks}[subsection]{Remarks}

\numberwithin{equation}{subsection}

% Macros
%
\def\lim{\mathop{\mathrm{lim}}\nolimits}
\def\colim{\mathop{\mathrm{colim}}\nolimits}
\def\Spec{\mathop{\mathrm{Spec}}}
\def\Hom{\mathop{\mathrm{Hom}}\nolimits}
\def\Ext{\mathop{\mathrm{Ext}}\nolimits}
\def\SheafHom{\mathop{\mathcal{H}\!\mathit{om}}\nolimits}
\def\SheafExt{\mathop{\mathcal{E}\!\mathit{xt}}\nolimits}
\def\Sch{\mathit{Sch}}
\def\Mor{\mathop{\mathrm{Mor}}\nolimits}
\def\Ob{\mathop{\mathrm{Ob}}\nolimits}
\def\Sh{\mathop{\mathit{Sh}}\nolimits}
\def\NL{\mathop{N\!L}\nolimits}
\def\CH{\mathop{\mathrm{CH}}\nolimits}
\def\proetale{{pro\text{-}\acute{e}tale}}
\def\etale{{\acute{e}tale}}
\def\QCoh{\mathit{QCoh}}
\def\Ker{\mathop{\mathrm{Ker}}}
\def\Im{\mathop{\mathrm{Im}}}
\def\Coker{\mathop{\mathrm{Coker}}}
\def\Coim{\mathop{\mathrm{Coim}}}

% Boxtimes
%
\DeclareMathSymbol{\boxtimes}{\mathbin}{AMSa}{"02}

%
% Macros for moduli stacks/spaces
%
\def\QCohstack{\mathcal{QC}\!\mathit{oh}}
\def\Cohstack{\mathcal{C}\!\mathit{oh}}
\def\Spacesstack{\mathcal{S}\!\mathit{paces}}
\def\Quotfunctor{\mathrm{Quot}}
\def\Hilbfunctor{\mathrm{Hilb}}
\def\Curvesstack{\mathcal{C}\!\mathit{urves}}
\def\Polarizedstack{\mathcal{P}\!\mathit{olarized}}
\def\Complexesstack{\mathcal{C}\!\mathit{omplexes}}
% \Pic is the operator that assigns to X its picard group, usage \Pic(X)
% \Picardstack_{X/B} denotes the Picard stack of X over B
% \Picardfunctor_{X/B} denotes the Picard functor of X over B
\def\Pic{\mathop{\mathrm{Pic}}\nolimits}
\def\Picardstack{\mathcal{P}\!\mathit{ic}}
\def\Picardfunctor{\mathrm{Pic}}
\def\Deformationcategory{\mathcal{D}\!\mathit{ef}}
\usepackage{cleveref}
\crefname{originallemma}{original lemma}{original lemmas}
\Crefname{originallemma}{Original lemma}{Original lemmas}
\makeatletter
\newlabel{categories-lemma-adjoint-exact}{{0038}{source}{}{sourcepointer-1}{}}
\newlabel{categories-lemma-adjoint-exact@cref}{{[originallemma][0][]0038}{[1][0][]source}{}{}{}}
\makeatother
\crefname{originalreference}{original reference}{original references}
\Crefname{originalreference}{Original reference}{Original references}
\makeatletter
\newlabel{lemma-connected-topological}{{0CAM}{source}{}{sourcepointer-2}{}}
\newlabel{lemma-connected-topological@cref}{{[originalreference][0][]0CAM}{[1][0][]source}{}{}{}}
\makeatother
\crefname{originalequation}{original equation}{original equations}
\Crefname{originalequation}{Original equation}{Original equations}
\makeatletter
\newlabel{modules-equation-towards-constructible-sets}{{0CAJ}{source}{}{sourcepointer-3}{}}
\newlabel{modules-equation-towards-constructible-sets@cref}{{[originalequation][0][]0CAJ}{[1][0][]source}{}{}{}}
\makeatother
\crefname{originallemma}{original lemma}{original lemmas}
\Crefname{originallemma}{Original lemma}{Original lemmas}
\makeatletter
\newlabel{modules-lemma-constructible-in-constant}{{0CAL}{source}{}{sourcepointer-4}{}}
\newlabel{modules-lemma-constructible-in-constant@cref}{{[originallemma][0][]0CAL}{[1][0][]source}{}{}{}}
\makeatother
\crefname{originallemma}{original lemma}{original lemmas}
\Crefname{originallemma}{Original lemma}{Original lemmas}
\makeatletter
\newlabel{modules-lemma-filtered-colimit-constructibles}{{0CAI}{source}{}{sourcepointer-5}{}}
\newlabel{modules-lemma-filtered-colimit-constructibles@cref}{{[originallemma][0][]0CAI}{[1][0][]source}{}{}{}}
\makeatother
\crefname{originallemma}{original lemma}{original lemmas}
\Crefname{originallemma}{Original lemma}{Original lemmas}
\makeatletter
\newlabel{sheaves-lemma-directed-colimits-sections}{{009F}{source}{}{sourcepointer-6}{}}
\newlabel{sheaves-lemma-directed-colimits-sections@cref}{{[originallemma][0][]009F}{[1][0][]source}{}{}{}}
\makeatother
\crefname{originallemma}{original lemma}{original lemmas}
\Crefname{originallemma}{Original lemma}{Original lemmas}
\makeatletter
\newlabel{sheaves-lemma-f-map}{{008K}{source}{}{sourcepointer-7}{}}
\newlabel{sheaves-lemma-f-map@cref}{{[originallemma][0][]008K}{[1][0][]source}{}{}{}}
\makeatother
\crefname{originaldefinition}{original definition}{original definitions}
\Crefname{originaldefinition}{Original definition}{Original definitions}
\makeatletter
\newlabel{topology-definition-spectral-space}{{08YG}{source}{}{sourcepointer-8}{}}
\newlabel{topology-definition-spectral-space@cref}{{[originaldefinition][0][]08YG}{[1][0][]source}{}{}{}}
\makeatother
\crefname{originallemma}{original lemma}{original lemmas}
\Crefname{originallemma}{Original lemma}{Original lemmas}
\makeatletter
\newlabel{topology-lemma-spectral-sub}{{0902}{source}{}{sourcepointer-9}{}}
\newlabel{topology-lemma-spectral-sub@cref}{{[originallemma][0][]0902}{[1][0][]source}{}{}{}}
\makeatother
\title{Sections of a sheaf on a spectral space agree with sections on a closed subset meeting every closure connectedly}
\author{The Stacks Project Authors}
\date{stacks-project:etale-cohomology.tex:local-snapshot-20261011; GNU FDL 1.2 or later}
\begin{document}
\maketitle
\noindent\textit{Editorial note.} All delivered mathematical text is the author's own, copied from the pinned Stacks Project chapter file etale-cohomology.tex; no statement or proof step was rewritten, reordered or removed.\par\medskip
\noindent\textit{Editorial note.} The counted proof is the author's complete argument: the proof is written out in full: the injectivity argument by specialisation, the two reductions (filtered colimits and the embedding), the inheritance of the hypotheses by the closed pieces, and the final constant-sheaf case.\par\medskip
\noindent\textit{Editorial note.} The label lemma-h0-topological names the counted statement, and the proof environment immediately following it in the source is delivered verbatim as its proof.\par\medskip
\noindent\textit{Editorial note.} The counted statement is carried by the author's own lemma environment in the chapter source.\par\medskip
\noindent\textit{Editorial note.} The statement and proof are copied from the pinned chapter file; the nine cited inputs are kept as external pointers with their canonical Stacks tags.\par\medskip
\noindent\textit{Allowed standard prerequisites (editorial).}\begin{itemize}
\item Spectral spaces and quasi-compact opens. Topology, the definition of a spectral space and the lemma that a closed subset with the induced properties is again spectral. Supplies the setting and the inheritance of the hypotheses by the closed subsets Z\_i.
\item Description of a sheaf as a filtered colimit of constructible-like sheaves. Modules, the equation describing the auxiliary sheaves and the lemma that every sheaf is a filtered colimit of them. Supplies the reduction to a sheaf of the auxiliary form via commutation with filtered colimits.
\item Embedding into a finite product of pushforwards of constant sheaves. Modules, the lemma embedding a sheaf of the auxiliary form into a finite product of (Z\_i -> X)\_* of constant sheaves. Supplies the reduction to the case of a finite constant sheaf.
\item The connectedness decomposition for spectral spaces. Same chapter, the lemma turning a decomposition of a closed subset with connected closures-intersections into a decomposition of the whole space. Supplies the final case of a finite constant sheaf.
\end{itemize}\medskip
\ifdefined\hypertarget\hypertarget{sourcepointer-1}{}\fi
\noindent\textit{Source reference.} Original lemma 0038: sources/stacks-chow/tags.txt, source lines 125--125. Prerequisite pointer invoked at proof line(s) 15127; it names a standard fact of a sibling chapter that the counted argument uses, and every step of the counted proof built on it is retained verbatim in the delivered body.\par\medskip
\ifdefined\hypertarget\hypertarget{sourcepointer-2}{}\fi
\noindent\textit{Source reference.} Original reference 0CAM: sources/stacks-etale-cohomology/pinned-original-source.tex, source lines 15072--15072. Prerequisite pointer invoked at proof line(s) 15162; it names a standard fact of the same chapter that the counted argument uses, and every step of the counted proof built on it is retained verbatim in the delivered body.\par\medskip
\ifdefined\hypertarget\hypertarget{sourcepointer-3}{}\fi
\noindent\textit{Source reference.} Original equation 0CAJ: sources/stacks-chow/tags.txt, source lines 15439--15439. Prerequisite pointer invoked at proof line(s) 15124; it names a standard fact of a sibling chapter that the counted argument uses, and every step of the counted proof built on it is retained verbatim in the delivered body.\par\medskip
\ifdefined\hypertarget\hypertarget{sourcepointer-4}{}\fi
\noindent\textit{Source reference.} Original lemma 0CAL: sources/stacks-chow/tags.txt, source lines 15441--15441. Prerequisite pointer invoked at proof line(s) 15147; it names a standard fact of a sibling chapter that the counted argument uses, and every step of the counted proof built on it is retained verbatim in the delivered body.\par\medskip
\ifdefined\hypertarget\hypertarget{sourcepointer-5}{}\fi
\noindent\textit{Source reference.} Original lemma 0CAI: sources/stacks-chow/tags.txt, source lines 15438--15438. Prerequisite pointer invoked at proof line(s) 15125; it names a standard fact of a sibling chapter that the counted argument uses, and every step of the counted proof built on it is retained verbatim in the delivered body.\par\medskip
\ifdefined\hypertarget\hypertarget{sourcepointer-6}{}\fi
\noindent\textit{Source reference.} Original lemma 009F: sources/stacks-chow/tags.txt, source lines 351--351. Prerequisite pointer invoked at proof line(s) 15129; it names a standard fact of a sibling chapter that the counted argument uses, and every step of the counted proof built on it is retained verbatim in the delivered body.\par\medskip
\ifdefined\hypertarget\hypertarget{sourcepointer-7}{}\fi
\noindent\textit{Source reference.} Original lemma 008K: sources/stacks-chow/tags.txt, source lines 320--320. Prerequisite pointer invoked at proof line(s) 15128; it names a standard fact of a sibling chapter that the counted argument uses, and every step of the counted proof built on it is retained verbatim in the delivered body.\par\medskip
\ifdefined\hypertarget\hypertarget{sourcepointer-8}{}\fi
\noindent\textit{Source reference.} Original definition 08YG: sources/stacks-chow/tags.txt, source lines 11339--11339. Prerequisite pointer invoked at proof line(s) 15104; it names a standard fact of a sibling chapter that the counted argument uses, and every step of the counted proof built on it is retained verbatim in the delivered body.\par\medskip
\ifdefined\hypertarget\hypertarget{sourcepointer-9}{}\fi
\noindent\textit{Source reference.} Original lemma 0902: sources/stacks-chow/tags.txt, source lines 11395--11395. Prerequisite pointer invoked at proof line(s) 15160; it names a standard fact of a sibling chapter that the counted argument uses, and every step of the counted proof built on it is retained verbatim in the delivered body.\par\medskip
\begin{lemma}
\label{lemma-h0-topological}
Let $Z \subset X$ be a closed subset of a topological space $X$.
Assume
\begin{enumerate}
\item $X$ is a spectral space
(Topology, Definition \ref{topology-definition-spectral-space}), and
\item for $x \in X$ the intersection $Z \cap \overline{\{x\}}$
is connected (in particular nonempty).
\end{enumerate}
Then for any sheaf $\mathcal{F}$ on $X$ we have
$\Gamma(X, \mathcal{F}) = \Gamma(Z, \mathcal{F}|_Z)$.
\end{lemma}

\begin{proof}
If $x \leadsto x'$ is a specialization of points, then there is a
canonical map $\mathcal{F}_{x'} \to \mathcal{F}_x$ compatible with
sections over opens and functorial in $\mathcal{F}$. Since every point
of $X$ specializes to a point of $Z$ it follows that
$\Gamma(X, \mathcal{F}) \to \Gamma(Z, \mathcal{F}|_Z)$ is injective.
The difficult part is to show that it is surjective.

\medskip\noindent
Denote $\mathcal{B}$ be the set of all quasi-compact opens of $X$.
Write $\mathcal{F}$ as a filtered colimit $\mathcal{F} = \colim \mathcal{F}_i$
where each $\mathcal{F}_i$ is as in
Modules, Equation (\ref{modules-equation-towards-constructible-sets}).
See Modules, Lemma \ref{modules-lemma-filtered-colimit-constructibles}.
Then $\mathcal{F}|_Z = \colim \mathcal{F}_i|_Z$ as restriction to $Z$
is a left adjoint (Categories, Lemma \ref{categories-lemma-adjoint-exact} and
Sheaves, Lemma \ref{sheaves-lemma-f-map}).
By Sheaves, Lemma \ref{sheaves-lemma-directed-colimits-sections}
the functors $\Gamma(X, -)$ and $\Gamma(Z, -)$ commute with filtered colimits.
Hence we may assume our sheaf $\mathcal{F}$ is as in
Modules, Equation (\ref{modules-equation-towards-constructible-sets}).

\medskip\noindent
Suppose that we have an embedding $\mathcal{F} \subset \mathcal{G}$.
Then we have
$$
\Gamma(X, \mathcal{F}) =
\Gamma(Z, \mathcal{F}|_Z) \cap \Gamma(X, \mathcal{G})
$$
where the intersection takes place in $\Gamma(Z, \mathcal{G}|_Z)$.
This follows from the first remark of the proof because we can check
whether a global section of $\mathcal{G}$ is in $\mathcal{F}$ by looking
at the stalks and because every point of $X$ specializes to a point of $Z$.

\medskip\noindent
By Modules, Lemma \ref{modules-lemma-constructible-in-constant}
there is an injection $\mathcal{F} \to \prod (Z_i \to X)_*\underline{S_i}$
where the product is finite, $Z_i \subset X$ is closed, and $S_i$ is finite.
Thus it suffices to prove surjectivity for the sheaves
$(Z_i \to X)_*\underline{S_i}$. Observe that
$$
\Gamma(X, (Z_i \to X)_*\underline{S_i}) = \Gamma(Z_i, \underline{S_i})
\quad\text{and}\quad
\Gamma(X, (Z_i \to X)_*\underline{S_i}|_Z) =
\Gamma(Z \cap Z_i, \underline{S_i})
$$
Moreover, conditions (1) and (2) are inherited by $Z_i$; this is clear
for (2) and follows from
Topology, Lemma \ref{topology-lemma-spectral-sub} for (1). Thus it
suffices to prove the lemma in the case of a (finite) constant sheaf.
This case is a restatement of Lemma \ref{lemma-connected-topological}
which finishes the proof.
\end{proof}

\end{document}
